% ============================================================
%  Part VII — Imperative, Konjunktiv, Passive
% ============================================================
\gpart{VII}{Mood \& Voice}

% ------------------------------------------------------------
\gsec{Imperative}

\begin{gtbl}{colspec={X[0.8,l] X[1.2,l] X[1.4,l]}}
  Addressing & How to form it & Example \\
  du   & stem, no pronoun \eng{-e optional} & \textbf{Komm}! \textbf{Geh}! \\
  ihr  & stem + \en{-t}, no pronoun & \textbf{Kommt}! \textbf{Geht}! \\
  Sie  & Infinitiv + \textbf{Sie}   & \textbf{Kommen Sie}! \\
  wir  & Infinitiv + \textbf{wir}   & \textbf{Gehen wir}! \eng{let's go} \\
\end{gtbl}

\tcap{Stem changes in the du-form}

\begin{gtbl}{colspec={X[0.9,l] X[1,l] X[1.1,l]}}
  Verb & du-Imperativ & Note \\
  geben \eng{e $\rightarrow$ i}  & \textbf{Gib}!   & keeps the change \\
  sehen \eng{e $\rightarrow$ ie} & \textbf{Sieh}!  & keeps the change \\
  fahren \eng{a $\rightarrow$ ä} & \textbf{Fahr}!  & \textbf{loses} the umlaut \\
  laufen & \textbf{Lauf}! & loses the umlaut \\
  arbeiten & \textbf{Arbeite}! & -d/-t stems keep the -e \\
\end{gtbl}

\tcap{Irregular and separable}

\begin{flattbl}{colspec={X[1,l] X[1,l] X[1,l]}}
  sein   & Sei! & Seid! / Seien Sie! \\
  haben  & Hab! & Habt! / Haben Sie! \\
  werden & Werde! & Werdet! / Werden Sie! \\
  aufstehen & Steh auf! & Steht auf! / Stehen Sie auf! \\
\end{flattbl}

\begin{gnote}
Add \textit{bitte}, \textit{doch} or \textit{mal} to take the edge off:
\textit{Komm doch mal her!} A bare imperative sounds blunt in German, just as it
does in English.
\end{gnote}

% ------------------------------------------------------------
\gsec{Konjunktiv II}

Used for the unreal: hypotheticals, wishes, and polite requests.

\tcap{The two ways to build it}

\begin{gtbl}{colspec={X[0.95,l] X[2.05,l]}}
  Method & When \\
  \textbf{würde} + Infinitiv & the default for almost every verb \\
  Real Konjunktiv II form & \textit{sein}, \textit{haben}, the modals, and a handful of common strong verbs \\
\end{gtbl}

\begin{gtbl}{colspec={X[1.25,l] X[1,c] X[1,c] X[1,c]}}
  Person & würde & hätte & wäre \\
  ich       & würde   & hätte   & wäre   \\
  du        & würdest & hättest & wärst  \\
  er/sie/es & würde   & hätte   & wäre   \\
  wir       & würden  & hätten  & wären  \\
  ihr       & würdet  & hättet  & wärt   \\
  sie/Sie   & würden  & hätten  & wären  \\
\end{gtbl}

\tcap{Forms worth knowing outright}

\begin{flattbl}{colspec={X[1,l] X[1,l] X[1,l] X[1,l]}}
  können $\rightarrow$ könnte & müssen $\rightarrow$ müsste & dürfen $\rightarrow$ dürfte & sollen $\rightarrow$ sollte \\
  wollen $\rightarrow$ wollte & mögen $\rightarrow$ möchte  & wissen $\rightarrow$ wüsste & gehen $\rightarrow$ ginge \\
  kommen $\rightarrow$ käme   & geben $\rightarrow$ gäbe    & tun $\rightarrow$ täte      & lassen $\rightarrow$ ließe \\
\end{flattbl}

\tcap{What it is used for}

\begin{flattbl}{colspec={X[0.85,l] X[2.15,l]}}
  Politeness  & \textbf{Könnten} Sie mir helfen? \textbullet{} Ich \textbf{hätte} gern einen Kaffee. \\
  Hypothesis  & Wenn ich Zeit \textbf{hätte}, \textbf{würde} ich kommen. \\
  Wish        & Ich \textbf{wünschte}, ich \textbf{wäre} dort. \\
  Advice      & An deiner Stelle \textbf{würde} ich warten. \\
  Past unreal & Wenn ich Zeit \textbf{gehabt hätte}, \textbf{wäre} ich gekommen. \\
\end{flattbl}

\begin{gnote}
There is only one past Konjunktiv II, built with \textit{hätte} or \textit{wäre}
plus the Partizip II. \textit{Würde} is never used to form the past.
\end{gnote}

% ------------------------------------------------------------
\gsec{Passive}

\tcap{Vorgangspassiv — werden \eng{something is being done}}

\begin{gtbl}{colspec={X[1,l] X[2,l]}}
  Tense & Example \\
  Präsens        & Das Haus \textbf{wird} gebaut. \\
  Präteritum     & Das Haus \textbf{wurde} gebaut. \\
  Perfekt        & Das Haus \textbf{ist} gebaut \textbf{worden}. \\
  Plusquamperfekt & Das Haus \textbf{war} gebaut \textbf{worden}. \\
  Futur I        & Das Haus \textbf{wird} gebaut \textbf{werden}. \\
  With a modal   & Das Haus \textbf{muss} gebaut \textbf{werden}. \\
\end{gtbl}

\begin{gnote}
In the passive perfect the participle of \textit{werden} loses its \textit{ge-}
and becomes \textbf{worden}, never \textit{geworden}. That single letter is the
most common passive mistake.
\end{gnote}

\tcap{Zustandspassiv — sein \eng{something is in a state}}

\begin{flattbl}{colspec={X[1,l] X[1.5,l]}}
  Das Fenster \textbf{wird} geöffnet. & it is being opened \eng{the action} \\
  Das Fenster \textbf{ist} geöffnet.  & it is open \eng{the result} \\
\end{flattbl}

\tcap{Naming the agent}

\begin{flattbl}{colspec={X[0.75,l] X[1,l] X[1.6,l]}}
  von + Dativ  & a person or agent & Der Brief wurde \textbf{von mir} geschrieben. \\
  durch + Akk. & a means or cause  & Die Stadt wurde \textbf{durch das Feuer} zerstört. \\
\end{flattbl}

\begin{gnote}
Turning active into passive: the accusative object becomes the nominative
subject, and the old subject either disappears or comes back with \textit{von}.
If the active verb has no accusative object, the passive has no subject and
takes a dummy \textit{es}: \textit{Es wird getanzt.}
\end{gnote}

\tcap{Ways to dodge the passive}

\begin{flattbl}{colspec={X[1,l] X[1.6,l]}}
  man + active   & \textbf{Man} baut das Haus. \\
  sich lassen    & Das \textbf{lässt sich} machen. \eng{that can be done} \\
  sein + zu + Inf. & Das \textbf{ist zu machen}. \\
\end{flattbl}
